\chapter{Experimental Protocol}\label{ch:protocol}

The protocol is the main methodological product of the report. It is designed so that a result cannot be produced by an under-tuned baseline, by an uncounted cost, by too few seeds, or by choosing the threshold after seeing the data. This chapter describes it in the order in which a result is produced: what is counted (\cref{sec:flops}), when a run has ``reached'' a target (\cref{sec:crossing}), how every arm is tuned (\cref{sec:tuning}), how arms are compared (\cref{sec:stats}), and how hypotheses are fixed in advance (\cref{sec:prereg}).

\section{Design principles}
\begin{enumerate}
\item \textbf{Compute-matched.} The unit of cost is the number of counted FLOPs, so that a method pays for every extra network, extra forward pass or auxiliary loss it uses. Step counts are reported next to FLOPs, never instead of them.
\item \textbf{Target-based.} The outcome is the cost to reach a fixed validation accuracy (0.97 and 0.975), not accuracy at a fixed budget. This is the style of metric used by recent benchmarks of training algorithms \citep{dahl2023benchmarking}.
\item \textbf{Equally tuned.} Every arm is tuned by the same budget of 25 random-search trials \citep{bergstra2012random}; the baseline is tuned more than once so that the effect of re-tuning is measured.
\item \textbf{Paired and replicated.} The same seeds are used in every group, so every comparison is a paired difference; the decisive comparisons use 20 seeds.
\item \textbf{Pre-registered.} Hypotheses, margins, intervals and decision rules are written down and locked in version control before the confirmatory runs.
\end{enumerate}

\section{Data and seeds}
The data split is fixed (55\,000 training / 5\,000 validation images, \cref{sec:setting}) and the test set is not used. Seeds are partitioned by purpose so that no seed is used for both tuning and evaluation: selection seeds for tuning are at least 100 (in practice 1\,000 or more) and never overlap the evaluation seeds; the Study~1 seeds are 0--19 (0--9 for secondary groups); Study~2 uses seeds 1300--1319; Study~3 uses seeds 2300--2319; exploratory looks use other ranges (listed in \cref{app:reproducibility}). A seed fixes the initialisation and the order of the training batches. Because the validation set has 5\,000 images, a single evaluation of a model with true accuracy $0.975$ has a sampling standard error of $\sqrt{0.975\cdot0.025/5000}\approx0.22$ percentage points. Differences between arms smaller than this cannot be resolved at the level of a single curve; the paired design and the number of seeds address this at the level of means.

\section{What is counted: FLOPs}\label{sec:flops}
FLOPs are counted with PyTorch's \code{FlopCounterMode}, which counts the matrix multiplications of the forward and backward passes (a multiply--add is two FLOPs). The count is cumulative and logged at every evaluation. It \emph{includes} the forward and backward passes of every network in an arm (for arm D the network $G$ and its losses; for arm S the full forward pass and the partial backward pass; for Selective-Backprop every selection forward pass). It does \emph{not} include elementwise operations (activations, softmax, loss), the optimiser update, the top-$k$ selection, the gather of rows, data loading, or evaluation passes. This choice is common and favours methods whose overhead is dominated by non-matrix operations, which is why wall-clock time is measured separately (\cref{sec:wallclock}) and why \cref{sec:limitations} lists FLOPs-versus-time as a limitation.

\section{When a run has reached a target}\label{sec:crossing}
Validation accuracy is evaluated every 50 steps. A run \emph{reaches} a target at the first evaluation $i$ at which the accuracy is at or above the target at both $i$ and $i+1$; the reported step and FLOPs are those of evaluation $i$. A run whose only evaluation at or above the target is the last one has not reached it. A run that never reaches the target is scored with the FLOPs of its full run (so a method that fails to reach the target is never rewarded) and counted in a column \code{n\_never}.

\begin{protocolbox}[title={A post-hoc protocol decision, stated openly}]
The first version of the rule scored the \emph{first} evaluation at or above the target. In the raw logs of an earlier tuning pass, 118 of the 182 runs that crossed $0.975$ were back below it at the next evaluation ($65\%$), so a single lucky evaluation would decide the winner of a tuning. The rule was therefore changed to two consecutive evaluations on 22 October, \emph{after} the first tuning results had been seen but \emph{before} the locking of the first pre-registration and before any final run. Re-scoring all runs under three rules showed that the rule changes which trial wins for the baseline, but not the standing of the meeting loss (ratio to A: $1.31$, $1.19$ and $1.25$ in-sample under first crossing, two consecutive and a 3-evaluation mean). The target $0.98$ is only secondary because the untuned baseline misses it in 3 of 5 seeds under the adopted rule.
\end{protocolbox}

\section{Tuning every arm the same way}\label{sec:tuning}
Each arm is tuned by random search with 25 trials. A trial is trained for the full 3\,000 steps on two selection seeds, and its score is the mean over these seeds of the FLOPs to reach $0.975$ under the crossing rule above (lower is better; ties are broken by the higher mean validation accuracy over the last five evaluations). The search space is the same for every arm for the shared hyper-parameters, plus the arm-specific weights (\cref{tab:space}). The winner is run on fresh seeds. Trials are sampled once per (tag, arm) and recorded, so that a tuning can be reproduced exactly.

\begin{table}[t]
\centering\small
\caption{Search space of the random search (25 trials per arm).}
\label{tab:space}
\begin{tabular}{@{}l l l@{}}
\toprule
Hyper-parameter & Distribution & Used by\\
\midrule
learning rate & log-uniform $[10^{-4},10^{-2}]$ & all arms\\
weight decay & log-uniform $[10^{-6},10^{-2}]$ & all arms\\
activation & \{ReLU, $\tanh$\} & all arms\\
learning-rate schedule & \{constant, cosine\} & all arms\\
$\lambda_{\mathrm{meet}}$ & log-uniform $[10^{-3},10]$ & D, V1--V3, F\\
$\lambda_{\mathrm{inv}}$ & log-uniform $[0.1,10]$ & D, V1, V2\\
warm-up steps $T_{\mathrm{w}}$ & \{200, 500, 1000\} & D, V1--V3, F\\
$\lambda_{\mathrm{aux}}$ & log-uniform $[10^{-2},1]$ & B, B-code, C\\
$\sigma_{\mathrm{inv}}$ & log-uniform $[0.01,0.3]$ & V1\\
lower-layer learning-rate multiplier & log-uniform $[0.5,8]$ & E\\
\bottomrule
\end{tabular}
\end{table}

\paragraph{The noise floor of tuning.} A tuning is itself a random procedure. To measure how much the answer depends on it, the baseline was tuned several times with fully independent sampled trials and selection seeds: the first tuning (A1, tag \code{t1}), a second one (A2, \code{t2}), and four more (A5--A8, \code{t5}--\code{t8}). The gain claimed for any new arm has to exceed the differences between these re-tunings (\cref{sec:stats}). The chosen configurations are in \cref{app:hyperparameters}. The result of this check is one of the main findings of the report (\cref{sec:noise}).

\section{Comparing arms}\label{sec:stats}
\paragraph{Statistic.} For two groups run on the same seeds, the statistic is the paired percentage change of the mean FLOPs-to-target of group $G$ relative to a reference group $R$:
\begin{equation}
\Delta = 100\cdot\left(\frac{\overline{c}_G}{\overline{c}_R}-1\right),
\end{equation}
where $\overline{c}$ is the mean over seeds of the cost. Negative values mean that $G$ is cheaper. An interval for $\Delta$ is obtained by the percentile bootstrap \citep{efron1993introduction}, resampling seeds jointly for both groups (10\,000 resamples, bootstrap seed 0).

\paragraph{Intervals and the error budget.} Each pre-registration tests a small family of hypotheses and splits an error budget of $0.05$ among them: two-sided $97.5\%$ intervals for the two comparisons of Studies~1 and~2, and $98.3\%$ ($=1-0.05/3$) intervals for the three comparisons of Study~3. No further correction is made for the two targets because both must agree.

\paragraph{Three-way verdict.} Each hypothesis states a margin $m$ in advance (a relative saving that is worth having, $15\%$ for the main comparisons). With point estimate $\hat\Delta$ and interval $[\mathrm{lo},\mathrm{hi}]$ at a target:
\begin{itemize}
\item \emph{supported} if $\hat\Delta\le -m$ and $\mathrm{hi}<0$;
\item \emph{refuted} if $\mathrm{lo}>-m$ (the interval excludes a saving of $m$ or more);
\item \emph{inconclusive} otherwise.
\end{itemize}
The overall verdict requires the two primary targets, $0.97$ and $0.975$, to agree. Because a ``null result'' has an explicit definition, a failure of an idea can be reported as such and not as a lack of significance. A simulation done before the first confirmatory runs, with a noise level estimated from earlier runs, shows what the rule can resolve: at a true change of $0$ the verdict is ``refuted'' in $96\%$ of simulations with 20 seeds ($79\%$ with 10); at a true saving of $20\%$ it is ``supported'' in $92\%$ ($83\%$); at a true saving of exactly the margin of $15\%$ it is ``supported'' about half of the time. A true gain just above the margin therefore often ends as ``inconclusive''.

\paragraph{Re-tuning shift.} For a claim against the baseline, a ``supported'' verdict additionally requires the gain to exceed the \emph{shift} caused by re-tuning the baseline. In Study~1 the shift is the absolute change of the second tuning A2 against A1 at the same target. In Studies~2 and~3 it is the \emph{largest reduction} in mean FLOPs-to-target against A1, on the same seeds, achieved by any of the five re-tunings A2, A5--A8 (zero if none is below A1). The latter rule is a maximum over five noisy estimates, which is biased upward and so makes a ``supported'' verdict harder, not easier. It was chosen after seeing that A2 alone differs from A1 by $+34\%$ on the Study~1 seeds (a worse re-tuning does not weaken a claim of the form ``S beats the baseline'', whereas the absolute rule would have made the verdict inconclusive by itself); this is recorded as a post-hoc decision in \cref{app:deviations}.

\section{Pre-registration, evidence tiers and the drift log}\label{sec:prereg}
Hypotheses, margins, decision rules, seeds, configurations and the analysis code of each confirmatory test are committed to the git repository and marked with a tag (\cref{tab:prereg}). The run script refuses to start unless the tag exists, is an ancestor of the current commit, and the locked paths (source code, tuned configurations, hypothesis and method documents, the plan) are unchanged since the tag. Any later change to these paths is a \emph{drift} and must be written down, with its reason, in a drift log; the report lists the relevant entries in \cref{app:deviations}. Only results of locked tests are \tierC{}; everything else is labelled \tierE{} and is allowed to generate hypotheses but not to support claims. After a locked test has been run, its seed set is not extended: an ``inconclusive'' result is not followed by adding seeds to the same test.

\begin{table}[t]
\centering\small
\caption{The three pre-registrations. A run header records the commit and whether the working tree was clean; all runs of the three confirmatory tests were started from the lock commit with a clean tree.}
\label{tab:prereg}
\begin{tabularx}{\textwidth}{@{}l l l >{\raggedright\arraybackslash}X >{\raggedright\arraybackslash}p{2.6cm} l@{}}
\toprule
Tag & Lock commit & Date & Hypotheses (margin) & Runs, seeds & Interval\\
\midrule
\code{prereg-v1} & \code{3378d36} & 23 Oct & H1 meeting loss vs.\ A (15\%); H6 deep supervision vs.\ A (15\%) & 150 runs, seeds 0--19 / 0--9 & 97.5\%\\
\code{prereg-v2} & \code{cc2830c} & 26 Oct & H7a S30 vs.\ A1 (15\%); H7b S30 vs.\ tuned batch-48 baseline (0\%) & 220 runs, seeds 1300--1319 & 97.5\%\\
\code{prereg-v3} & \code{a49e9e3} & 30 Oct & H8 S20w0 vs.\ A1 (15\%); H9 vs.\ S30 (10\%); H10 vs.\ tuned Selective-Backprop (15\%) & 260 runs, seeds 2300--2319 & 98.3\%\\
\bottomrule
\end{tabularx}
\end{table}

\section{Computing environment}
All runs were executed on one laptop CPU (WSL2 on x86-64), PyTorch 2.14.1 (CPU build) and Python 3.14.4 (as recorded for the timing experiment), one torch thread per process and four worker processes in parallel. A 3\,000-step run takes about a minute of effective time under this load; the 260 runs of Study~3 took about 4.4 hours. Since the runs do not use a GPU and the model is small, none of the compute is a barrier to repeating the experiments (\cref{app:reproducibility}).
